\section{Prior work}
\label{sec:prior}

The idea that epicycles and Fourier-type expansions are the same thing, and
the individual ingredients of our construction, have a long history. We
organize it by the ingredient it supplies. The survey is the result of a
literature search whose limits we state plainly: we cannot exclude that
relevant older or non-English sources escaped it, and several primary sources
(Bessel's 1824 memoir, Schiaparelli's 1874 work) were not consulted directly
and are cited through secondary accounts or not at all.

\paragraph*{Epicycles as a universal representation.}
Hanson \cite{Hanson1960} argued that epicyclic astronomy is mathematically
powerful enough to reproduce essentially any observed path, and that its
decline was not due to weak mathematics. This is the classical statement of
the ``epicycles can draw anything'' thesis, repeated in popular treatments
\cite{Mathologer2018,3b1b2019,WikiDeferent} and in a short approximation-theory
note \cite{Acosta2019}. It is a kinematic statement and is not derived from
dynamics. Its philosophical consequence, that a system able to fit any
curve is hard to refute, is discussed by Carman \cite{Carman2010,Carman2015}.
Manin \cite{Manin2018} and Boccaletti \cite{Boccaletti2001} place the
continuity of circle-composition ideas in broader historical context.

\paragraph*{Epicycles as quasi-periodic motion.}
The closest conceptual precedent is Gallavotti \cite{Gallavotti2001}. He
rewrites a quasi-periodic Fourier series as a deferent carrying nested
epicycles ordered by amplitude, expands the Kepler ellipse of Mars in powers
of the eccentricity to obtain a deferent and two epicycles at first order and
more at higher order, and argues, through Kolmogorov's theorem, that motions
near integrable ones are again epicyclic, with approximation only over finite
times in general. Our work does not supersede this; as far as we can tell,
it is planar, does not write the geocentric difference of two orbits
explicitly, and does not use Bessel coefficients or treat three dimensions.
The adjacent mathematical framework is Bohr's theory of almost-periodic
functions \cite{Bohr1925}, which is where we place the geocentric
expansion. Scholz \cite{Scholz2022} shows that moving from heliocentric
circular orbits to an Earth-based observer leads to epicyclic constructions,
which is our geocentric transform restricted to circular orbits.

\paragraph*{The Bessel--Fourier expansion of Kepler motion.}
That a Kepler orbit has an exact Fourier expansion in the mean anomaly with
Bessel-function coefficients is classical. It originates with Lagrange and
Bessel \cite{Dutka1995,Dutka1997} and its convergence theory is that of
Kapteyn series \cite{Kapteyn1893,Watson1944}; expansions of the radius and
anomalies are standard textbook material
\cite{BrouwerClemence1961,MurrayDermott1999}, surveyed by Colwell
\cite{Colwell1993}, with recent work on summation and convergence
\cite{Tautz2011,Borghi2024}. A point worth keeping separate is that the
Laplace limit $e\approx 0.6627$ governs the power-series solution in $e$, not
the Fourier series in the mean anomaly, which converges for all $e<1$
\cite{Dutka1997,Colwell1993}. Our expansion rests on these results;
what we add is their assembly into a three-dimensional geocentric epicycle
chain with explicit truncation bounds (Sec.~\ref{sec:theory}).

\paragraph*{Ptolemaic models as truncations of Kepler motion.}
Evans showed that the equant reproduces Kepler's second law to first order in
eccentricity \cite{Evans1984} and traced the successive historical
approximations to Kepler motion for Mars \cite{Evans1988}. Rushkin
\cite{Rushkin2015} analyzed the accuracy attainable with Ptolemy's
eccentricity-division rule, finding errors of second order in $e$. Carman and
Recio \cite{CarmanRecio2019} give the most careful modern quantitative
comparison of Ptolemaic and Keplerian longitude models, with optimized
parameters. We do not claim to give the first Ptolemy-versus-Kepler
comparison. The standard historical sources for the models themselves are
\cite{Toomer1998,Neugebauer1975,SwerdlowNeugebauer1984,Evans1998}.

\paragraph*{Newtonian dynamics to epicycles.}
Several works derive epicycles from Newtonian orbits. Peralta, Calles and
Y\'epez \cite{Peralta2003} integrate Newton's equations and apply an FFT to the
resulting time series, showing that the Fourier terms generate the deferent
and epicycles of the Copernican model; this is the most direct precedent for
the numerical half of our programme, though heliocentric and planar, without
perturbations, and without the analytic Bessel transform. Sugon, Bragais and
McNamara \cite{Sugon2008} obtain Copernican epicycles from the linearized
equation of motion about a circular orbit, to first order and in a
heliocentric frame. Pucacco \cite{Pucacco2012} expands the Kepler problem in
a co-rotating frame, obtaining at first order a circle plus a second circle
at twice the frequency, which he calls an epicyclic expansion, and extends it
with normal forms to perturbed cases. Paillusson and
Booth \cite{Paillusson2020} compute geocentric transforms of heliocentric
trajectories to discuss the underdetermination of world systems, without a
Fourier decomposition. An unreviewed software repository \cite{epa058}
reconstructs planar geocentric trajectories from Keplerian elements as sums
of rotating circles using FFT peaks and least squares, which shows that the
geocentric numerical idea is in circulation, so our novelty cannot rest on
the idea alone.

\paragraph*{Frequency analysis and ephemeris series.}
Operationally, planetary positions are already represented as sums of uniform
circular motions. VSOP87 \cite{Bretagnon1988} expresses heliocentric
coordinates as long sums of Poisson terms, and each term with no secular
factor is a uniform circular motion; a geocentric version is obtained by
differencing two such theories. Kudryavtsev \cite{Kudryavtsev2016} obtained
analytical series for DE431 by spectral fitting of a numerical ephemeris.
Laskar's numerical frequency analysis \cite{Laskar1990,Laskar1992,Laskar1993,%
Laskar1999} recovers quasi-periodic decompositions with an accuracy far
better than a plain FFT, with convergence estimates in terms of the data
length \cite{Laskar2003}, and also marks the limits of the epicycle picture:
the motion of the solar system is chaotic, and over long times
frequencies drift so that a fixed set of epicycles ceases to represent the
solution \cite{Laskar1990,FuLaskar2019}. We use NAFF as our numerical
tool and reuse these limits in our statement of validity.

\paragraph*{Position of the present work.}
To summarize, concept and framing are due to Gallavotti
\cite{Gallavotti2001}; the Bessel expansion to Lagrange, Bessel and
Kapteyn; the FFT demonstration to Peralta et al.\
\cite{Peralta2003}; first-order analytic derivations to Sugon et al.\
\cite{Sugon2008} and Pucacco \cite{Pucacco2012}; and the practical
representation to VSOP87 and Laskar-type frequency analysis. We found no
source that combines (i) the explicit closed-form three-dimensional
geocentric expansion with ellipse-as-two-circles structure, (ii) rigorous
truncation bounds, (iii) a frequency-analysis study that checks NAFF against the analytic
expansion on unperturbed geocentric trajectories and, for a perturbed
(N-body) geocentric Ceres trajectory, compares its reconstruction error with
fitted-Kepler baselines, with the limits of that comparison stated, and (iv) a precise statement of the domain of validity. That
combination, together with the visualization of
Sec.~\ref{sec:viz}, is what this paper contributes, as an expository
synthesis with new quantitative results rather than as a new discovery.
